% Lab report: aim, theory, apparatus, procedure, data, analysis with
% uncertainties, and conclusion.
\documentclass[11pt,a4paper]{article}

\usepackage[T1]{fontenc}
\usepackage[margin=2.5cm]{geometry}
\usepackage{amsmath}
\usepackage{graphicx}
\usepackage{booktabs}
\usepackage{siunitx}     % numbers with units: \qty{9.81}{\metre\per\second\squared}
\usepackage[hidelinks]{hyperref}

\title{Experiment 1: Measuring $g$ with a Simple Pendulum}
\author{Your Name \\ Lab partner: Partner Name \\ Course code, Section}
\date{Performed: \today}

\begin{document}

\maketitle

\section{Aim}
To determine the acceleration due to gravity from the period of a simple
pendulum.

\section{Theory}
For small oscillations a pendulum of length $L$ has period
\begin{equation}
  T = 2\pi\sqrt{\frac{L}{g}},
  \qquad\text{so}\qquad
  g = \frac{4\pi^{2}L}{T^{2}}.
  \label{eq:g}
\end{equation}

\section{Apparatus}
Pendulum bob, string, retort stand, metre rule (least count \qty{1}{\milli\metre}),
stopwatch (least count \qty{0.01}{\second}).

\section{Procedure}
\begin{enumerate}
  \item Measure the length $L$ from the pivot to the centre of the bob.
  \item Release the bob from a small angle and time 20 oscillations.
  \item Repeat for five lengths.
\end{enumerate}

\section{Observations}

\begin{table}[htbp]
  \centering
  \caption{Time for 20 oscillations at each length.}
  \label{tab:data}
  \begin{tabular}{S[table-format=1.3] S[table-format=2.2] S[table-format=1.3] S[table-format=1.2]}
    \toprule
    {$L$ (\unit{\metre})} & {$t_{20}$ (\unit{\second})} & {$T$ (\unit{\second})} & {$T^{2}$ (\unit{\second\squared})} \\
    \midrule
    0.400 & 25.40 & 1.270 & 1.61 \\
    0.500 & 28.38 & 1.419 & 2.01 \\
    0.600 & 31.08 & 1.554 & 2.41 \\
    0.700 & 33.58 & 1.679 & 2.82 \\
    0.800 & 35.90 & 1.795 & 3.22 \\
    \bottomrule
  \end{tabular}
\end{table}

\section{Analysis}
A plot of $T^{2}$ against $L$ is a straight line of slope $4\pi^{2}/g$.
From Table~\ref{tab:data} the slope is \qty{4.03}{\second\squared\per\metre}, giving
\[
  g = \frac{4\pi^{2}}{4.03} = \qty{9.80}{\metre\per\second\squared}.
\]

\subsection{Uncertainty}
With relative uncertainties $\delta L / L$ and $\delta T / T$,
\[
  \frac{\delta g}{g} = \sqrt{\left(\frac{\delta L}{L}\right)^{2} + \left(2\,\frac{\delta T}{T}\right)^{2}}.
\]

\section{Result}
$g = \qty{9.80(5)}{\metre\per\second\squared}$, in agreement with the accepted value.

\section{Sources of Error}
\begin{itemize}
  \item Reaction time when starting and stopping the stopwatch.
  \item The small-angle approximation in Eq.~\eqref{eq:g}.
\end{itemize}

\end{document}
